% Generated draft: requires mathematical and source review.
\documentclass[11pt]{article}
\usepackage[margin=25mm]{geometry}
\usepackage{fontspec}
\setmainfont{Latin Modern Roman}
\usepackage{amsmath,amssymb,mathtools,unicode-math}
\setmathfont{Latin Modern Math}
\usepackage{enumitem,needspace,xurl,hyperref}
\setlength{\parindent}{0pt}
\setlength{\parskip}{5pt}
\setlength{\emergencystretch}{3em}
\newcommand{\N}{\mathbb{N}}
\newcommand{\Z}{\mathbb{Z}}
\newcommand{\Q}{\mathbb{Q}}
\newcommand{\R}{\mathbb{R}}
\newcommand{\C}{\mathbb{C}}
\newcommand{\F}{\mathbb{F}}
\newcommand{\A}{\mathbb{A}}
\newcommand{\PP}{\mathbb{P}}
\newcommand{\E}{\mathbb{E}}
\newcommand{\eps}{\varepsilon}
\newcommand{\dd}{\,\mathrm{d}}
\begin{document}
\section*{UnsolvedMath ID: 700629}
\textbf{Source problem code:} catalogue entry
\subsection{Definitions and mathematical statement}
A semigroup \(S\) has cancellation if, for all \(x,y,z\in S\),
\[
xy=xz\Longrightarrow y=z,\qquad yx=zx\Longrightarrow y=z.
\]
Let \(A\subseteq S\) be a finite semigroup generating set. Write \(S^1\) for \(S\) with an identity adjoined if necessary, and define
\[
B_{S,A}(n)=\{a_1\cdots a_k:0\leq k\leq n,\ a_i\in A\}\subseteq S^1,
\qquad b_{S,A}(n)=|B_{S,A}(n)|,
\]
where the empty product is the identity. Interpret the source's term ``non-exponential growth'' in the standard subexponential sense:
\[
\forall c>1\ \exists n_0\in\N\ \forall n\geq n_0,
\qquad b_{S,A}(n)<c^n.
\]
The source does not explicitly define this growth terminology.

A group of left quotients of \(S\), denoted \(G=S^{-1}S\), means a group \(G\) together with an injective semigroup homomorphism \(\iota:S\to G\) such that
\[
G=\{\iota(s)^{-1}\iota(t):s,t\in S\}.
\]
Put \(X=\iota(A)\cup\iota(A)^{-1}\). This is a finite group generating set for \(G\). Define
\[
B_{G,X}(n)=\{x_1\cdots x_k:0\leq k\leq n,\ x_i\in X\},
\qquad b_{G,X}(n)=|B_{G,X}(n)|.
\]
Exponential growth means
\[
\exists c>1\ \exists n_0\in\N\ \forall n\geq n_0,
\qquad b_{G,X}(n)\geq c^n.
\]
These growth properties are independent of the respective finite generating sets.

\textbf{Question.} Does there exist a finitely generated cancellative semigroup \(S\) of subexponential growth whose group of left quotients \(G=S^{-1}S\) has exponential growth?

The source explicitly states that the group of left quotients exists under its hypotheses; the question above retains this requirement and does not ask whether such a group exists for an arbitrary cancellative semigroup. The source also states that an affirmative answer would give a positive solution to Problem~12 in S.~Wagon's \emph{The Banach--Tarski Paradox} (1985). The text of that problem is not supplied, so this consequence is retained as a source attribution without further formulation or verification.
\subsection{Short English statement}
R.~I.~Grigorchuk asks whether a finitely generated semigroup with both cancellation laws can have subexponential growth while its group of left quotients has exponential growth. The source attributes to an affirmative answer a positive solution of Problem~12 in Wagon's 1985 book; that auxiliary problem is not supplied.
\subsection{Sources}
\begin{itemize}
\item Supplied local UnsolvedMath registry record: permanent ID \(700629\), exact source code \texttt{catalogue entry}, title \emph{Kourovka Notebook Problem 10.12}. The registry labels the problem ``open''; this is supplied status metadata, not an independently verified assessment.
\item \emph{The Kourovka Notebook (Unsolved Problems in Group Theory)}, Problem~10.12, proposed by R.~I.~Grigorchuk in the 10th issue (1986). The supplied background identifies the 21st issue edition (2026), arXiv identifier \texttt{1401.0300v48}. The supplied collection description names E.~I.~Khukhro and V.~D.~Mazurov as editors. These bibliographic details have not been independently checked.
\item S.~Wagon, \emph{The Banach--Tarski Paradox}, Cambridge University Press, 1985, Problem~12. Cited in the supplied statement for the claimed consequence of an affirmative answer; the auxiliary problem and the implication have not been independently checked.
\end{itemize}
\end{document}
